\chapter{讨论与局限}

\section{讨论}

本项目没有另起一套 \texttt{bensz-postdoc} 公共包，而是选择在现有 \texttt{bensz-thesis} 产品线内注册独立的 \texttt{thesis-smu-postdoc} template/profile/style。这样做有两个好处：第一，继续复用已经稳定、已被多套项目验证的构建链路和页面样式能力；第二，既保证博士后研究报告在业务和打包层面独立成立，又避免为高度相近的底层排版能力重复维护第二套公共包。

\section{当前公开示例的边界}

本模板仍然是公开演示项目，因此存在明确边界：

\begin{enumerate}
  \item 不包含任何真实患者数据、私有统计结果或未公开研究发现；
  \item 不承诺已经对齐某一年度所有院内细则，仅优先吸收全国统一规则与封面核心字段；
  \item 不主动引入 Word 到 DOCX 的双向导出链路，当前交付重点是稳定 PDF。
\end{enumerate}

\section{适用建议}

实际使用时，建议优先替换以下内容：

\begin{enumerate}
  \item \texttt{extraTex/meta.tex} 中的封面与题名页元信息；
  \item \texttt{extraTex/front/abstract\_*.tex} 中的摘要内容；
  \item \texttt{extraTex/body/} 中各章正文；
  \item \texttt{extraTex/back/} 中的附录、简历、成果和通信地址。
\end{enumerate}
